\section{Related Work}
\label{sec:related}

\subsection{Classical spoofing detection}
\label{sec:rw_classical}

Classical anti-spoofing tests operate on quantities the tracking loops already produce,
which makes them inexpensive to compute and straightforward to interpret. Received power
monitoring compares the carrier-to-noise density ratio $C/N_0$ against the value expected
for a satellite at a given elevation; a counterfeit transmitted above the authentic power
raises this quantity, and the additional energy also registers in the automatic gain
control setting of the front end \citep{akos2012gnss}. Correlation shape tests examine the
peak formed when the incoming signal is correlated against a locally generated replica of
the spreading code. An authentic signal yields a symmetric peak whose width is set by the
chip duration and the front end bandwidth, whereas a counterfeit present alongside the
authentic signal superimposes a second peak and distorts the composite, producing
asymmetry that can be quantified from early, prompt and late correlator taps
\citep{psiaki2016gnss}. Consistency tests exploit relations that must hold among
measurements of a genuine signal: the code and carrier observables of one satellite track
the same changing geometric range and must therefore evolve together, so a counterfeit
that displaces one without the other violates the relation \citep{psiaki2016gnss}.

How much of the adversary's parameter space these tests actually cover has itself been
measured, and the answer is not all of it. The detection rate of signal quality monitoring
metrics falls for particular combinations of counterfeit code phase and carrier phase, and
metrics formed from several correlator offsets have been proposed to widen that coverage,
reaching a reported 96.1 percent across those combinations \citep{wu2025musd}. Reading the
correlator output as an eye diagram is reported to increase coverage by more than 60
percent over signal quality monitoring and carrier-to-noise density tests
\citep{wu2025eyediagram}. What a classical test achieves therefore depends on where in its
parameter space the spoofer chooses to operate. The present paper asks the same question of
learned detectors, and asks it of transmissions that satisfy the conditions for capturing a
receiver.

These tests are complementary in the sense that no one of them subsumes another, and a
receiver running several raises the cost of a successful attack. They nonetheless share a
structural property. Each is defined by a single named physical signature, which is what
makes it interpretable, and a spoofer that maintains the named quantity within its nominal
range satisfies the test. Because the family of deployed tests is documented in the open
literature, an adversary can treat the signature set as known and design a transmission
that respects it. This observation motivates detection rules that are not specified in
advance by a human.

\subsection{Machine learning based spoofing detection}
\label{sec:rw_learning}

A learned detector replaces a hand specified statistic with a decision rule inferred from
labelled data. A receiver is run over recordings in which the ground truth is known, each
epoch yields a feature vector of the observables available at that instant together with a
label, and a classifier is fitted to reproduce those labels before being applied to unseen
epochs. Because the classifier consumes the observables jointly, it can respond to
configurations that no univariate statistic encodes.

Two model families recur in this literature, and both are represented in the present
study. Single epoch classifiers operate on one observation vector at a time. Decision
trees partition the feature space by recursive threshold tests, and ensembles such as
random forests and gradient boosted trees aggregate many such partitions to reduce
variance. Nearest neighbour classifiers assign a label from the closest training examples,
and support vector machines construct a maximum margin separating surface, in the
nonlinear case through a kernel. Sequence models instead consume a window of consecutive
epochs and can therefore exploit temporal structure. Convolutional networks extract local
patterns across the window, recurrent architectures such as the long short term memory
propagate state along it \citep{hochreiter1997long}, and attention based architectures
weight the epochs of the window against one another \citep{vaswani2017attention}.

Applications to GNSS spoofing span both families. Deep networks have been applied to
correlator outputs \citep{borhanidarian2024deep}, to signal quality features on small
uncrewed aircraft \citep{sung2022cnn}, to receiver exchange format records
\citep{romaniuc2025lstm} and to unmanned aerial platforms \citep{yang2025deep}. Classical
models have been applied through support vector machines \citep{chen2022svm}, nearest
neighbour classifiers \citep{mahroof2024ml} and feed forward networks benchmarked against
recurrent alternatives \citep{jullian2022deep}. Real time operation has been addressed
\citep{li2025realtime}, as has enlargement of the training distribution by generative
augmentation \citep{chen2025cgan}. Reported detection accuracy across this body of work is
consistently high, and the architectures achieving it are heterogeneous, which suggests
that the reported figures are substantially determined by the evaluation corpora and
not by architectural choices.

\subsection{Generalization beyond the training attacks}
\label{sec:rw_generalization}

A classifier fitted to recorded examples characterises the attacks contained in those
recordings, and whether its accuracy persists on an attack outside them is a separate
question that has received direct attention.

One response treats detection as one class learning, fitting only the authentic
distribution so that any sufficiently atypical observation is flagged, which removes the
dependence on a catalogue of attack examples \citep{iqbal2024deep}. A second augments the
training set with synthetically generated attacks in order to widen the distribution the
detector has seen \citep{chen2025cgan}. A third constrains the detector with physics.
Code carrier consistency, the same relation the classical consistency test evaluates, has
been imposed as a training constraint, so the network must respect a relation authentic
signals satisfy. This is reported to preserve accuracy across held out scenarios on which
unconstrained detectors degrade \citep{song2026generalized}.
The appeal of the third approach is that a physically grounded invariant does not depend
on which attacks were recorded.

This line of work addresses distribution shift, in which the test attacks are drawn from a
different distribution than the training attacks but are still drawn and not chosen.
Adversarial selection is a distinct axis: the attack is not sampled from any distribution
but optimised against the deployed system. The two questions are complementary, and a
library of transmitted attacks of the kind constructed here provides material against
which invariant based detectors can subsequently be examined.

\subsection{Adversarial robustness of learned classifiers}
\label{sec:rw_adversarial}

Classifiers fitted to data can be driven across their own decision boundaries by small,
deliberately chosen input perturbations, an observation that originated in image
classification and has since been shown to be general
\citep{szegedy2014intriguing,goodfellow2015explaining}. The associated literature covers
attack construction under model access \citep{madry2018towards}, methodological standards
for evaluating defensive claims \citep{carlini2019evaluating}, and the framing of the
problem as a security question \citep{biggio2018wild}.

Two results from this literature bear directly on the evaluation of spoofing detectors.
The first concerns adversary knowledge. Attacks that differentiate the loss require access
to model parameters, whereas decision based attacks require only the classifier's output
label and are nonetheless effective \citep{brendel2018decision,chen2020hopskipjump}. A
model that withstands a gradient attack has therefore not been shown robust, because a
model that supplies uninformative gradients may still be crossed by a search that queries
only decisions, a failure mode termed gradient masking \citep{athalye2018obfuscated}. The
second concerns cost. Apparent robustness obtained by displacing a decision threshold is
not a gain, since any classifier can be made difficult to cross by flagging nearly all
inputs; the accuracy and robustness of a fitted classifier trade against one another
\citep{tsipras2019robustness}. A robustness claim is therefore interpretable only when
reported jointly with the false alarm rate that purchased it, and the present study adopts
that convention throughout.

\subsection{Physical realizability and the evaluation gap}
\label{sec:rw_physical}

Transferring these results to a sensing system introduces a requirement absent from the
image setting. A perturbation expressed as a displacement of a feature vector corresponds
to an attack only if some physical action on the sensor produces that displacement.
Research on other sensing modalities has addressed this by constructing attacks in the
physical domain, including against automotive ranging sensors
\citep{petit2015remote,cao2021invisible}, against speech recognition
\citep{carlini2018audio} and against synthetic aperture radar imagery \citep{lin2023sar}.
In the GNSS setting, adversarial perturbation of detector inputs has been reported
\citep{an2025adversarial}.

For a GNSS receiver the requirement is stringent, because the observables a detector reads
are not independent coordinates. Several are computed from the same correlator
accumulations within a single integration period, so displacing one while holding the
others fixed describes a receiver state that no incident waveform produces. Establishing
that a given feature space perturbation is realizable therefore requires an argument about
receiver processing that is specific to the observable set in use. Evaluating at the
waveform level removes the need for such an argument by construction, since the adversary
acts on transmitter parameters and the observables are whatever the receiver computes while
tracking the resulting signal. This is the approach adopted in the present study, and
Section~\ref{sec:res_compare} reports the corresponding feature space evaluation alongside
it so that the two can be compared directly.
